\section{总结}
本文围绕智能城市生成系统的软件设计展开说明，从系统概述、架构分解、数据组织、面向对象建模和界面设计五个层面给出了较为完整的设计描述。报告明确了主系统、Blender 插件执行端、解析服务、仿真服务以及数据与资源存储之间的职责边界，也对项目、场景版本、任务、结构化命令、插件作业、执行结果和日志等核心对象进行了统一组织。

从体系结构上看，系统采用任务驱动的组织方式：用户输入首先在主系统中完成接收、解释和确认，再由任务编排机制驱动插件执行与结果回写；从数据设计上看，系统以项目、版本和任务为主线，以仿真对象和操作日志作为辅助支撑，从而满足结果回看、版本比较和过程追溯的需要；从面向对象设计上看，报告通过用例图、类图、活动图、顺序图、协作图、状态图、构件图和部署图对系统的静态结构、动态行为、软件构成和运行环境进行了系统化表达；从界面设计上看，系统以统一登录入口和分角色工作台为交互基础，使建模、分析和治理三类活动能够在一致的导航逻辑下完成。

综合而言，本报告所描述的软件设计方案体现了三个核心特征：其一，以结构化任务作为自然语言输入与几何生成执行之间的中间桥梁，从而控制解析结果的可解释性与可执行性；其二，以场景版本和任务状态管理为基础，实现生成、修改、分析和归档过程的连续衔接；其三，以角色分工和构件分层为原则，使交互层、控制层、执行层和存储层之间保持清晰边界。上述设计为系统后续实现、接口细化、模块协同和运行部署提供了统一依据。
